\chapter{Casos de uso}\label{cap:casos}

\section{Visão geral}
Um caso de uso descreve o objetivo do ator e o caminho até um resultado.
Precondições delimitam o estado inicial; alternativas são resultados previstos;
exceções interrompem a operação sem atingir a pós-condição de sucesso.
As etapas separam ações do ator e respostas do sistema.

\begin{quadrosrs}
\caption{Atores e objetivos dos casos de uso}
\begin{tabelasrs}{L{0.15\textwidth}YY}
\cabecalhosrs\textbf{Caso} & \textbf{Ator e objetivo} & \textbf{Requisitos}\\\midrule
\ucref{UC-001} & Participante: inscrever-se. & \reqref{RF-002}, \reqref{RF-003}.\\
\ucref{UC-002} & Participante: cancelar inscrição. & \reqref{RF-004}.\\
\ucref{UC-003} & Visitante: consultar eventos. & \reqref{RF-001}.\\
\ucref{UC-004} & Participante: consultar seus registros. & \reqref{RF-005}.\\
\end{tabelasrs}
\fonte{Elaboração didática deste guia.}
\end{quadrosrs}

\Needspace{0.50\textheight}
\section{UC-001: realizar inscrição}\label{uc:UC-001}
\begin{fichasrs}{Condições de UC-001}
\campoficha{Ator}{Participante autenticado (STK-002).}
\campoficha{Objetivo}{Obter inscrição confirmada para um evento.}
\campoficha{Precondições}{Sessão válida; evento conhecido. Vaga, publicação e prazo são verificados durante o fluxo, não presumidos pela interface.}
\campoficha{Gatilho}{Participante confirma a intenção de se inscrever.}
\campoficha{Pós-condição}{Uma inscrição ativa do titular é persistida e seu identificador é apresentado; somente uma vaga é consumida.}
\campoficha{Relações}{\reqref{RF-002}, \reqref{RF-003}; \rnref{RN-001}, \rnref{RN-002}; \ctref{CT-002}, \ctref{CT-003}, \ctref{CT-006}.}
\end{fichasrs}

\begin{quadrosrs}
\caption{Fluxo principal de UC-001}
\begin{tabelasrs}{L{0.06\textwidth}YY}
\cabecalhosrs\textbf{Nº} & \textbf{Ação do ator} & \textbf{Resposta do sistema}\\\midrule
1 & Seleciona um evento do catálogo. & Mostra evento, início e vagas para confirmação.\\
2 & Confirma a inscrição. & Valida sessão e obtém a identidade do titular.\\
3 & Aguarda o resultado. & Verifica inscrição ativa existente; se não existe, valida publicação, prazo e vaga.\\
4 & --- & Persiste a inscrição e o efeito na capacidade de forma atômica.\\
5 & Consulta a confirmação. & Retorna o identificador persistido e mostra o resultado.\\
\end{tabelasrs}
\fonte{Elaboração didática deste guia.}
\end{quadrosrs}

\begin{quadrosrs}
\caption{Alternativas e exceções de UC-001}
\begin{tabelasrs}{L{0.16\textwidth}L{0.13\textwidth}Y}
\cabecalhosrs\textbf{Caminho} & \textbf{Etapa} & \textbf{Comportamento e término}\\\midrule
A1: duplicidade & 3 & Retornar inscrição ativa existente, inclusive após o prazo; encerrar sem nova criação.\\
A2: falta de vaga & 3 & Informar falta de vagas; encerrar sem gravar.\\
A3: inelegibilidade & 3 & Negar evento não publicado ou prazo encerrado; encerrar sem gravar.\\
E1: sessão inválida & 2 & Solicitar autenticação; não ler nem alterar inscrição pessoal.\\
E2: falha de gravação & 4 & Informar indisponibilidade; não confirmar nem consumir vaga.\\
E3: resposta perdida & 5 & Repetição encontra o registro confirmado e segue A1, sem duplicar.\\
\end{tabelasrs}
\fonte{Elaboração didática deste guia.}
\end{quadrosrs}

\Needspace{0.50\textheight}
\section{UC-002: cancelar inscrição}\label{uc:UC-002}
\begin{fichasrs}{Condições de UC-002}
\campoficha{Ator}{Titular autenticado de uma inscrição (STK-002).}
\campoficha{Objetivo}{Desistir da inscrição e liberar sua vaga.}
\campoficha{Precondições}{Sessão válida; identificador de inscrição conhecido pelo participante. O servidor verifica titularidade e prazo.}
\campoficha{Gatilho}{Participante confirma o cancelamento.}
\campoficha{Pós-condição}{Registro cancelado, instante do primeiro cancelamento registrado e uma única vaga liberada.}
\campoficha{Relações}{\reqref{RF-004}; \rnref{RN-001}, \rnref{RN-002}; \ctref{CT-004}, \ctref{CT-006}, \ctref{CT-009}.}
\end{fichasrs}

\begin{quadrosrs}
\caption{Fluxo principal de UC-002}
\begin{tabelasrs}{L{0.06\textwidth}YY}
\cabecalhosrs\textbf{Nº} & \textbf{Ação do ator} & \textbf{Resposta do sistema}\\\midrule
1 & Abre suas inscrições e seleciona um registro ativo. & Mostra evento e opção de cancelar quando elegível.\\
2 & Confirma o cancelamento. & Valida sessão, titularidade e instante anterior ao início.\\
3 & --- & Altera a situação e libera uma vaga na mesma operação persistida.\\
4 & Consulta o resultado. & Confirma o cancelamento e atualiza a situação exibida.\\
\end{tabelasrs}
\fonte{Elaboração didática deste guia.}
\end{quadrosrs}

\begin{quadrosrs}
\caption{Alternativas e exceções de UC-002}
\begin{tabelasrs}{L{0.16\textwidth}L{0.13\textwidth}Y}
\cabecalhosrs\textbf{Caminho} & \textbf{Etapa} & \textbf{Comportamento e término}\\\midrule
A1: já cancelada & 2 & Se ainda no prazo, confirmar sem alterar o instante original ou as vagas.\\
A2: prazo encerrado & 2 & Negar no início ou depois dele; conservar o registro.\\
E1: sessão inválida & 2 & Solicitar autenticação; não alterar dados.\\
E2: ausente/não autorizada & 2 & Responder como recurso inexistente, sem expor dados de outra pessoa.\\
E3: falha de gravação & 3 & Informar indisponibilidade; conservar situação e vagas anteriores.\\
\end{tabelasrs}
\fonte{Elaboração didática deste guia.}
\end{quadrosrs}

\Needspace{0.50\textheight}
\section{UC-003: consultar eventos}\label{uc:UC-003}
\begin{fichasrs}{Especificação de UC-003}
\campoficha{Ator e objetivo}{Visitante (STK-001) deseja encontrar um evento futuro.}
\campoficha{Precondições}{Base de eventos local disponível; sessão não é necessária.}
\campoficha{Gatilho}{Visitante abre o catálogo.}
\campoficha{Fluxo principal}{1. Visitante solicita a lista. 2. Sistema seleciona eventos publicados e futuros, calcula vagas e ordena. 3. Interface apresenta nome, início e vagas.}
\campoficha{Alternativa A1}{Na seleção, nenhum evento é elegível: retornar lista vazia e informar ausência de resultados.}
\campoficha{Exceção E1}{Na leitura, armazenamento indisponível: informar falha temporária e permitir nova tentativa.}
\campoficha{Pós-condição}{Eventos elegíveis apresentados, sem alteração de dados.}
\campoficha{Relações}{\reqref{RF-001}, \reqref{RNF-001}; \ctref{CT-001}, \ctref{CT-005}.}
\end{fichasrs}

\Needspace{0.50\textheight}
\section{UC-004: consultar próprias inscrições}\label{uc:UC-004}
\begin{fichasrs}{Especificação de UC-004}
\campoficha{Ator e objetivo}{Participante (STK-002) deseja acompanhar suas inscrições.}
\campoficha{Precondições}{Sessão válida.}
\campoficha{Gatilho}{Participante abre a área de inscrições pessoais.}
\campoficha{Fluxo principal}{1. Participante solicita a lista. 2. Sistema obtém a identidade da sessão e seleciona somente seus registros. 3. Interface mostra os dados em ordem decrescente de criação e oferece cancelamento nos registros elegíveis.}
\campoficha{Alternativa A1}{Na seleção, não há registros do titular: apresentar lista vazia, sem incluir inscrições de outra pessoa.}
\campoficha{Exceção E1}{Na validação, sessão ausente ou expirada: solicitar autenticação, sem retornar dados pessoais.}
\campoficha{Exceção E2}{Na leitura, armazenamento indisponível: informar falha temporária.}
\campoficha{Pós-condição}{Registros do titular apresentados, sem mudança de situação ou vagas.}
\campoficha{Relações}{\reqref{RF-005}, \reqref{RNF-002}; \ctref{CT-006}, \ctref{CT-008}.}
\end{fichasrs}
